%%%%%%%%%%%%%%%%%%%%%%%%%%%%%%%%%%%%%%%%%%%%%%%%%%%%%%%%%%%%%%%%%%%%
%% I, the copyright holder of this work, release this work into the
%% public domain. This applies worldwide. In some countries this may
%% not be legally possible; if so: I grant anyone the right to use
%% this work for any purpose, without any conditions, unless such
%% conditions are required by law.
%%%%%%%%%%%%%%%%%%%%%%%%%%%%%%%%%%%%%%%%%%%%%%%%%%%%%%%%%%%%%%%%%%%%

\documentclass[
  digital,     %% The `digital` option enables the default options for the
               %% digital version of a document. Replace with `printed`
               %% to enable the default options for the printed version
               %% of a document.
%%  color,       %% Uncomment these lines (by removing the %% at the
%%               %% beginning) to use color in the printed version of your
%%               %% document
  oneside,     %% The `oneside` option enables one-sided typesetting,
               %% which is preferred if you are only going to submit a
               %% digital version of your thesis. Replace with `twoside`
               %% for double-sided typesetting if you are planning to
               %% also print your thesis. For double-sided typesetting,
               %% use at least 120 g/m² paper to prevent show-through.
  nosansbold,  %% The `nosansbold` option prevents the use of the
               %% sans-serif type face for bold text. Replace with
               %% `sansbold` to use sans-serif type face for bold text.
  nocolorbold, %% The `nocolorbold` option disables the usage of the
               %% blue color for bold text, instead using black. Replace
               %% with `colorbold` to use blue for bold text.
  lof,         %% The `lof` option prints the List of Figures. Replace
               %% with `nolof` to hide the List of Figures.
  lot,         %% The `lot` option prints the List of Tables. Replace
               %% with `nolot` to hide the List of Tables.
]{fithesis4}
%% The following section sets up the locales used in the thesis.
\usepackage[resetfonts]{cmap} %% We need to load the T2A font encoding
\usepackage[T1,T2A]{fontenc}  %% to use the Cyrillic fonts with Russian texts.
\usepackage[
  main=english, %% By using `czech` or `slovak` as the main locale
                %% instead of `english`, you can typeset the thesis
                %% in either Czech or Slovak, respectively.
  english, german, czech, slovak %% The additional keys allow
]{babel}        %% foreign texts to be typeset as follows:
%%
%%   \begin{otherlanguage}{german}  ... \end{otherlanguage}
%%   \begin{otherlanguage}{czech}   ... \end{otherlanguage}
%%   \begin{otherlanguage}{slovak}  ... \end{otherlanguage}
%%
%%
%% The following section sets up the metadata of the thesis.
\thesissetup{
    date        = 2026/02/24,
    university  = mu,
    faculty     = fi,
    type        = bc,
    department  = ?,
    author      = Adam David Malý,
    gender      = m,
    advisor     = {Honza + Martin},
    title       = {Implementation of MCP Server for Red Hat OpenShift AI Workbenches},
    TeXtitle    = {Implementation of MCP Server for Red Hat OpenShift AI Workbenches},
    keywords    = {MCP, RHOAI, Workbench, Server, Implementation},
    TeXkeywords = {MCP, RHOAI, Workbench, Server, Implementation},
    abstract    = {
      This thesis presents MCP server implementation for RHOAI workbenches.
    },
}
\usepackage{makeidx}      %% The `makeidx` package contains
\makeindex                %% helper commands for index typesetting.
%% These additional packages are used within the document:
\usepackage{paralist} %% Compact list environments
\usepackage{amsmath}  %% Mathematics
\usepackage{amsthm}
\usepackage{amsfonts}
\usepackage{url}      %% Hyperlinks
\usepackage{markdown} %% Lightweight markup
\usepackage{listings} %% Source code highlighting
\lstset{
  basicstyle      = \ttfamily,
  identifierstyle = \color{black},
  keywordstyle    = \color{blue},
  keywordstyle    = {[2]\color{cyan}},
  keywordstyle    = {[3]\color{olive}},
  stringstyle     = \color{teal},
  commentstyle    = \itshape\color{magenta},
  breaklines      = true,
}
\usepackage{floatrow} %% Putting captions above tables
\floatsetup[table]{capposition=top}
\usepackage[babel]{csquotes} %% Context-sensitive quotation marks
\usepackage[backend=biber, style=numeric]{biblatex}
\addbibresource{example.bib}
\begin{document}

%% ---------------------------------------------------------------------------
%% Chapter 1 -- Introduction
%% ---------------------------------------------------------------------------
\chapter{Introduction}
\label{chap:introduction}

Large language models (LLMs) have rapidly evolved in recent years from
conversational assistants into capable agents that can reason, plan, and take
actions in the real world. This shift has been made possible by the ability
to give LLMs access to external \textit{tools} -- functions that let a model
query a database, call an API, or manage infrastructure, all in response to a
natural language request. As the need for such integrations grows, a key
challenge emerges: how do different AI hosts and models connect to the
ever-expanding set of tools and data sources in a consistent way?

Anthropic introduced the Model Context Protocol (MCP) in November
2024~\parencite{mcp-announcement} to address exactly this problem. MCP is
an open standard that defines how AI applications connect to external tools
and data sources through a client-server architecture. Instead of each AI
host maintaining its own integrations, any MCP-compatible host
(Cursor, Claude Desktop, or a custom agent) can connect to any MCP server
and immediately gain access to the capabilities it exposes. This has already
led to a growing ecosystem of MCP servers for services ranging from databases
and file systems to cloud APIs and developer tools.

Red Hat OpenShift AI (RHOAI) is an enterprise machine learning platform
built on top of Red Hat OpenShift (which is based on Kubernetes). One of its
core features are workbenches -- containerised JupyterLab environments that
data scientists and engineers use for model training and experimentation.
Managing clusters with many workbench instances can be a tedious and
error-prone task, currently requiring navigation of the RHOAI web UI or
direct manipulation of Kubernetes resources. The goal of this thesis is to
eliminate that problem by allowing users to manage their workbenches entirely
through natural language -- simply describing what they want to an AI agent,
which then carries out the necessary operations on their behalf.

This thesis presents the design and implementation of an MCP server for
RHOAI workbenches, written in Go using the official MCP Go
SDK~\parencite{mcp-go-sdk}. The server exposes a comprehensive set of tools
covering the full lifecycle of workbench management: listing, creating,
updating, and deleting workbenches; managing notebook images and hardware
profiles; handling persistent volume claims; querying namespaces and pods;
and monitoring resource consumption at the workbench, namespace, user, and
cluster level. A safety-oriented read-only mode is provided by default, with
write operations gated behind an explicit environment variable. The server
also exposes MCP resources and a guided prompt to assist agents in
constructing correct workbench creation requests. Quality is validated
through unit tests and automated evaluation using the Promptfoo
framework~\parencite{promptfoo}, which measures tool selection accuracy,
parameter extraction correctness, and execution success rate.


%% ---------------------------------------------------------------------------
%% Chapter 2 -- Background: AI & MCP
%% ---------------------------------------------------------------------------
\chapter{Background: Artificial Intelligence and the Model Context Protocol}
\label{chap:ai-mcp}

This chapter introduces large language models and the agentic paradigm of tool
use, then describes the Model Context Protocol in detail.

\section{Large Language Models and AI Assistants}
\label{sec:llms}

At the core of modern AI assistants are \textit{large language models}
(LLMs) -- neural networks trained to predict and generate text by learning
statistical patterns from vast amounts of data. The dominant architectural
foundation for LLMs is the \textit{Transformer}, introduced by Vaswani et
al.\ in 2017~\parencite{vaswani2017attention}. The Transformer's
self-attention mechanism allows the model to relate every token in a
sequence to every other token, enabling it to capture long-range
dependencies in text far more effectively than earlier recurrent
architectures. Another benefit of the Transformer is that it is parallelizable,
making it possible to train large models on extensive datasets in a reasonable
time.

Modern LLMs are trained in two main stages. In the first stage,
\textit{pre-training}, the model is exposed to hundreds of billions of
tokens of text and learns to predict the next token in a sequence. This
process, while deceptively simple, causes the model to internalise broad
world knowledge, reasoning patterns, and linguistic structure. In the second
stage, the pre-trained model is fine-tuned to behave as a helpful assistant,
typically through a combination of supervised instruction fine-tuning and
reinforcement learning from human feedback
(RLHF)~\parencite{ouyang2022rlhf}. The result is a model that responds to
natural-language instructions in a coherent and useful way.

The products of this pipeline -- systems such as OpenAI's ChatGPT, Anthropic's
Claude, and Google's Gemini -- are commonly referred to as \textit{AI
assistants}. An AI assistant takes a user's message as input and produces a
textual response. It is fundamentally \textit{stateless} and
\textit{passive}: it answers questions, drafts documents, and explains
concepts, but it does not take actions in the world, does not maintain
persistent memory across sessions by default, and cannot access information
beyond its training data or the content of the current conversation. This
passivity is the key limitation that the agentic paradigm, described in the
next section, overcomes.

\section{AI Agents and Tool Use}
\label{sec:ai-agents}

Despite their impressive capabilities, pure language models are isolated from
the external world -- their knowledge does not extend beyond their training
data, they cannot query other systems, and they cannot take actions. The
concept of an \textit{AI agent} overcomes these limitations by equipping a
language model with \textit{tools} -- callable functions that the model can
invoke to interact with external systems, retrieve fresh information, or
modify state. Agents can also be given \textit{memory} -- the ability to
retain information across steps or sessions. This can range from simply
keeping the full conversation history in the model's context window, to
reading and writing records in an external database via dedicated memory
tools.

The key mechanism enabling this is \textit{tool calling} (also referred to
as \textit{function calling}). When a model is given a set of tool
definitions alongside a user request, it can -- instead of producing a plain
text response -- output a structured call to one of those tools, specifying
the tool name and the required arguments. The host application intercepts
this call, executes the corresponding function, and feeds the result back
into the model's context, after which the model continues its reasoning.
This loop can repeat multiple times within a single interaction, allowing the
model to chain tool calls to complete complex, multi-step tasks.

Yao et al.\ formalised this pattern in \textit{ReAct}~\parencite{yao2022react},
which unifies \textit{reasoning} and \textit{acting} into a single loop:
the model first produces a chain-of-thought reasoning trace -- thinking
through what it knows and what it needs -- and then selects and executes an
action (a tool call). The output returned by the tool becomes part of the
context for the next reasoning step. This structured loop gives agents a
degree of \textit{autonomy}: given a high-level goal, an agent can decompose
it into sub-tasks, select appropriate tools, and adapt its behaviour based on
the results it receives.

The practical impact of tool use is significant. An agent equipped with the
right tools can list resources in a Kubernetes cluster, query resource
consumption metrics, and delete workbench instances -- all in response to a
single natural-language instruction such as \enquote{delete workbenches
running for more than a week}. The Model Context Protocol provides a
standardised way of exposing those capabilities to any AI agent.

\section{The Model Context Protocol}
\label{sec:mcp}

\subsection{Protocol Overview}
\label{subsec:mcp-overview}

The Model Context Protocol (MCP) is an open standard introduced by Anthropic
in November 2024~\parencite{mcp-announcement} that defines how AI
applications connect to external tools and data sources. Its core motivation
is the \textit{M$\times$N integration problem}: without a shared standard,
connecting $M$ AI hosts to $N$ external services requires up to $M \times N$
custom integrations, each with its own authentication, serialisation, and
error-handling logic. MCP collapses this to $M + N$ by providing a single
protocol that both sides implement once.

The design of MCP draws inspiration from the Language Server Protocol
(LSP)~\parencite{lsp-spec}, which solved an analogous problem in software
development by standardising communication between code editors and language
analysis tools. Just as LSP allowed any compliant editor to work with any
compliant language server without bespoke glue code, MCP allows any
compliant AI host to work with any compliant tool server. The protocol
is transport-agnostic and message-based, built on top of
JSON-RPC~2.0~\parencite{jsonrpc-spec}. Official SDKs are available for
Python, TypeScript, Go, and other languages~\parencite{mcp-go-sdk}.

\subsection{MCP Clients and Servers}
\label{subsec:mcp-clients-servers}

The MCP architecture defines three distinct roles: \textit{hosts},
\textit{clients}, and \textit{servers}.

A \textbf{host} is the AI application that the end user interacts with
directly -- for example Claude Desktop, Cursor, or a custom agent
application. The host is responsible for managing the user interaction,
sending messages to the LLM, and orchestrating tool calls.

A \textbf{client} is a component that lives inside the host and manages
exactly one connection to one MCP server. A single host can run multiple
clients simultaneously, each connected to a different server, giving the
agent access to multiple sets of tools at once. The client speaks the MCP
protocol on behalf of the host: it sends requests to the server, receives
responses and notifications, and forwards results back to the LLM context.

A \textbf{server} is a lightweight process that exposes capabilities to
clients. A server typically wraps one specific external system -- a database,
a cloud API, a file system, or in the case of this thesis, Red Hat OpenShift
AI. The server does not communicate with the LLM directly; it only
communicates with the client that is connected to it. This separation of
concerns means that server authors do not need to know anything about which
LLM or host will ultimately consume their tools.

\subsection{Transport Mechanisms}
\label{subsec:mcp-transport}

MCP is transport-agnostic at the protocol level but defines two standard
transport mechanisms~\parencite{mcp-spec}.

\textbf{Standard I/O (stdio)} is the simplest transport. The host launches
the MCP server as a child process and communicates with it over the process's
standard input and standard output streams. Each JSON-RPC message is a
newline-delimited JSON object. This transport requires no network
configuration, has no authentication overhead, and is inherently scoped to
the local machine, making it well-suited for locally-installed tools. The
server implemented in this thesis uses the stdio transport, as reflected in
the entry point:

\begin{lstlisting}[language=Go]
server.Run(context.Background(), &mcp.StdioTransport{})
\end{lstlisting}

\textbf{HTTP with Server-Sent Events (SSE)} is the remote transport. The
client sends requests to the server via HTTP POST, and the server delivers
responses and unsolicited notifications back via an SSE stream. This allows
MCP servers to be deployed as network services, enabling multi-user scenarios
and remote infrastructure management. A newer \textit{Streamable HTTP}
transport, which consolidates both directions into a single HTTP endpoint, is
also defined in the specification as the recommended transport for new remote
deployments~\parencite{mcp-spec}.

\subsection{Tools, Resources, and Prompts}
\label{subsec:mcp-primitives}

MCP servers expose capabilities to clients through three primitive types:
tools, resources, and prompts~\parencite{mcp-spec}.

\textbf{Tools} are the primary mechanism for giving an AI agent the ability
to act. Each tool has a name, a human-readable description, and a JSON Schema
that describes its input parameters. When the LLM decides that a tool is
needed to fulfil a request, it emits a structured tool call containing the
tool name and arguments; the host executes the call via the client and
returns the result to the model's context. Tools are \textit{model-initiated}
-- it is the LLM that decides when and how to invoke them. This thesis
implements tools for the full workbench management lifecycle: listing,
creating, updating, and deleting workbenches, notebook images, hardware
profiles, persistent volume claims, and namespaces, as well as tools for
monitoring resource consumption.

\textbf{Resources} are URI-addressed data that the host can read and inject
into the model's context. Unlike tools, resources are \textit{host-initiated}:
it is the host or the user that decides which resources to include, not the
model. Resources are useful for providing structured reference data that the
model should be aware of when making decisions. This thesis exposes two
resources: an \textit{Image Catalog}, which lists all notebook images
available in the cluster with their versions and Python dependencies; and a
\textit{Default Hardware Profile}, which describes the default CPU and memory
allocation for new workbenches. An agent can read these resources before
calling the Create Workbench tool to ensure it selects a valid image and
reasonable resource limits.

\textbf{Prompts} are pre-defined, parameterisable message templates that
guide the model's behaviour for a specific task. A prompt can be invoked by
the user or the host and expands into one or more messages that are injected
into the conversation. This thesis provides a \textit{Create Workbench}
prompt that instructs the agent to first consult the Image Catalog resource
to select a suitable image and then invoke the Create Workbench tool -- a
pattern that combines all three MCP primitives in a single guided workflow.

%% ---------------------------------------------------------------------------
%% Chapter 2 -- Background: RHOAI & Workbenches
%% ---------------------------------------------------------------------------
\chapter{Background: Red Hat OpenShift AI and Workbenches}
\label{chap:rhoai}

\section{Kubernetes and OpenShift}
\label{sec:kubernetes}

\section{Red Hat OpenShift AI}
\label{sec:rhoai}

\section{Workbenches}
\label{sec:workbenches}

\section{MCP Tooling in RHOAI}
\label{sec:mcp-in-rhoai}

%% ---------------------------------------------------------------------------
%% Chapter 3 -- Analysis and Design
%% ---------------------------------------------------------------------------
\chapter{Analysis and Design}
\label{chap:analysis}

\section{Motivation and Goals}
\label{sec:motivation}

\section{Requirements}
\label{sec:requirements}

\section{Architecture}
\label{sec:architecture}

\section{Technology Choices}
\label{sec:technologies}

%% ---------------------------------------------------------------------------
%% Chapter 4 -- Implementation
%% ---------------------------------------------------------------------------
\chapter{Implementation}
\label{chap:implementation}

\section{Project Structure}
\label{sec:project-structure}

\section{MCP Tools}
\label{sec:tools}

\subsection{Workbench Management}
\label{subsec:workbench-tools}

\subsection{Hardware Profiles}
\label{subsec:hardware-tools}

\subsection{Notebook Images}
\label{subsec:image-tools}

\subsection{Storage and Namespaces}
\label{subsec:storage-tools}

\subsection{Resource Consumption}
\label{subsec:resource-tools}

\section{MCP Resources and Prompts}
\label{sec:resources-prompts}

\section{Security}
\label{sec:security}

\section{Code Quality}
\label{sec:code-quality}

\section{Evaluation}
\label{sec:evaluation}

%% ---------------------------------------------------------------------------
%% Chapter 5 -- Conclusion
%% ---------------------------------------------------------------------------
\chapter{Conclusion}
\label{chap:conclusion}

\section{Future Work: Notebooks 2.0}
\label{sec:future-work}

%% ---------------------------------------------------------------------------
\printbibliography[heading=bibintoc]

\end{document}
